\chapter{Nivel 4: Plataforma Central ThingsBoard y Apache Kafka}
\label{chap:server-thingsboard}

% ============================================================================
% PROPÓSITO ÚNICO: Detallar infraestructura de ingesta masiva y lógica de
% negocio en servidor central. Se agota el tema ThingsBoard Server aquí.
% ============================================================================

\section{Introducción}
\label{sec:server-intro}

% TODO: Introducción al capítulo (200-300 palabras)
% - Rol del servidor central: Gestión masiva dispositivos
% - ThingsBoard como plataforma IoT multi-tenant
% - Apache Kafka para streaming de alta concurrencia
% - Estructura del capítulo

\section{Infraestructura de Ingesta Masiva de Datos}
\label{sec:server-infraestructura}

% TODO: MIGRAR desde 04Implementacion_NEW.tex §4.5 (si existe)
% + EXPANDIR con arquitectura Kafka (1500 palabras NUEVAS)
% - Hardware: PC workstation 16 cores, 32 GB RAM
% - Kafka: Manejo concurrencia multi-dispositivo
% Longitud estimada: 6-8 páginas

\subsection{Especificaciones del Hardware}
\label{sec:server-hardware}

% TODO: NUEVO contenido (500 palabras)
% - CPU: AMD Ryzen 9 5950X (16 cores @ 3.4 GHz)
% - RAM: 32 GB DDR4-3200
% - Storage: NVMe 1 TB (PostgreSQL data)
% - Network: 10 GbE NIC (WAN connectivity)
% Tabla: Specs servidor vs alternativas cloud (AWS EC2)

\subsection{Arquitectura Apache Kafka}
\label{sec:server-kafka}

% TODO: NUEVO contenido (1500 palabras) - CRÍTICO
% Este es el contenido más NUEVO del capítulo
% - Kafka: Distributed streaming platform
% - Componentes:
%   * Kafka broker (3 replicas, factor 2)
%   * Zookeeper ensemble (3 nodes)
%   * Schema Registry (Avro schemas)
% - Topics:
%   * tb.mqtt.telemetry (partition 8, retention 7 days)
%   * tb.rule-engine.main (partition 4)
%   * tb.notifications (partition 2)
% Longitud crítica: 4-5 páginas

\subsubsection{Topología de Topics y Particiones}
\label{sec:server-kafka-topics}

% TODO: NUEVO contenido (800 palabras)
% - Topic tb.mqtt.telemetry:
%   * 8 particiones (hash deviceId)
%   * Replication factor: 2
%   * Retention: 7 días (168 horas)
%   * Compression: LZ4
% - Consumer groups:
%   * tb-rule-engine-consumer (parallel 8)
%   * tb-telemetry-processor (parallel 4)
% Figura: Kafka topic architecture

\subsubsection{Configuración de Producción}
\label{sec:server-kafka-config}

% TODO: NUEVO contenido (700 palabras)
% - Broker config:
%   * log.retention.hours=168
%   * num.partitions=8
%   * compression.type=lz4
%   * min.insync.replicas=1
% - Producer config (TB Edge):
%   * acks=1 (leader ack only)
%   * batch.size=16384
%   * linger.ms=10 (batching)
% Código: server.properties Kafka

\subsection{Balanceo de Carga y Rate Limiting}
\label{sec:server-balanceo}

% TODO: NUEVO contenido (600 palabras)
% - HAProxy: Load balancer MQTT
%   * Algoritmo: least-conn
%   * Health check: TCP port 1883
% - Rate limiting:
%   * 100 msgs/sec per device
%   * 10K msgs/sec total
% - Kafka throttling: produce.purgatory.purge.interval.requests
% Figura: Arquitectura load balancing

\section{Despliegue de ThingsBoard Server}
\label{sec:server-tb-server}

% TODO: MIGRAR desde 04Implementacion_NEW.tex §4.5
% - ThingsBoard CE 3.6.2
% - Docker Compose: 12 microservicios
% - Base de datos: PostgreSQL 14 + TimescaleDB 2.13
% Longitud estimada: 6-8 páginas

\subsection{Arquitectura de Microservicios}
\label{sec:server-tb-arquitectura}

% TODO: MIGRAR y EXPANDIR (600 palabras NUEVAS)
% - Microservicios ThingsBoard:
%   1. tb-core (8 instances, stateless)
%   2. tb-rule-engine (4 instances)
%   3. tb-web-ui (Nginx, 2 instances)
%   4. tb-js-executor (Node.js, 4 instances)
%   5. tb-mqtt-transport (2 instances)
%   6. tb-http-transport (2 instances)
%   7. tb-coap-transport (1 instance)
%   8. postgresql (primary + replica)
%   9. redis (cache + session)
%   10. kafka (3 brokers)
%   11. zookeeper (3 nodes)
%   12. prometheus + grafana (monitoring)
% Figura: Arquitectura microservicios ThingsBoard

\subsection{Configuración Docker Compose}
\label{sec:server-tb-docker}

% TODO: MIGRAR docker-compose.yml producción
% - Networks: tb-network (overlay)
% - Volumes: postgres-data, kafka-data, redis-data
% - Resource limits por servicio
% Código: docker-compose.yml servidor (>300 líneas)

\subsection{Optimización de PostgreSQL + TimescaleDB}
\label{sec:server-postgresql}

% TODO: NUEVO contenido (1000 palabras) - CRÍTICO
% - PostgreSQL 14: Relational data (devices, assets, users)
% - TimescaleDB 2.13: Time-series telemetry
%   * Hypertable: ts_kv (time-series key-value)
%   * Chunk interval: 7 días
%   * Retention policy: 90 días telemetry, 365 días aggregates
% - Índices optimizados:
%   * B-tree: entity_id, key
%   * BRIN: ts (timestamp)
% - Partitioning: Hash por entity_id (8 partitions)
% Longitud crítica: 3-4 páginas

\subsubsection{Diseño del Schema de Base de Datos}
\label{sec:server-postgresql-schema}

% TODO: NUEVO contenido (600 palabras)
% - Tablas principales:
%   * device (id, name, type, customer_id)
%   * ts_kv (entity_id, key, ts, bool_v, str_v, long_v, dbl_v)
%   * asset (id, name, type, customer_id)
%   * relation (from_id, to_id, relation_type)
% Figura: Diagrama ER schema ThingsBoard

\subsubsection{TimescaleDB Hypertables}
\label{sec:server-postgresql-hypertables}

% TODO: NUEVO contenido (800 palabras)
% - Hypertable ts_kv:
%   * Chunk interval: 7 días (168 horas)
%   * Compression: AFTER 30 days
%   * Retention: DROP AFTER 90 days
% - Continuous Aggregates:
%   * ts_kv_hourly (AVG, MAX, MIN)
%   * ts_kv_daily (SUM energy_kwh)
% Código: SQL creación hypertable + continuous aggregate

\subsection{Estrategia de Backup y Recuperación}
\label{sec:server-backup}

% TODO: NUEVO contenido (500 palabras)
% - Backup diario: pg_dump comprimido
% - Backup incremental: WAL archiving
% - Disaster recovery: RTO 4 horas, RPO 1 hora
% - Almacenamiento: S3-compatible (MinIO local)
% Código: Script backup.sh

\section{Lógica de Procesamiento de Negocio}
\label{sec:server-logica}

% TODO: MIGRAR desde 04Implementacion_NEW.tex §4.5
% + EXPANDIR Rule Chains (1200 palabras NUEVAS)
% - Asset Model: Tenant → Customer → Device
% - Rule Chains: Procesamiento telemetría
% - Alarm Management: Severidad CRITICAL/MAJOR/MINOR
% Longitud estimada: 7-9 páginas

\subsection{Modelado de Activos y Dispositivos}
\label{sec:server-assets}

% TODO: MIGRAR y EXPANDIR (800 palabras)
% - Asset Hierarchy:
%   * Tenant: Utility company (root)
%   * Customer: Building (192 devices)
%   * Asset: Building floors (3 floors × 64 devices)
%   * Device: ESP32-Thread-Meter (leaf)
% - Device Profile:
%   * Telemetry: 7 keys (voltage_v, current_a, etc.)
%   * Attributes: tariff_usd_kwh, location, install_date
%   * Alarms: 5 types (voltage_sag, over_current, etc.)
% Figura: Árbol de Assets ThingsBoard

\subsection{Diseño de Rule Chains}
\label{sec:server-rule-chains}

% TODO: NUEVO contenido (1500 palabras) - CRÍTICO
% Este apartado documenta lógica de negocio central
% - Rule Chains implementadas:
%   1. Root Rule Chain (default)
%   2. Energy Billing Chain
%   3. Alarm Processing Chain
%   4. Analytics Chain
% Longitud crítica: 4-5 páginas

\subsubsection{Root Rule Chain}
\label{sec:server-rule-chain-root}

% TODO: NUEVO contenido (600 palabras)
% - Nodes:
%   1. Message Type Switch (POST_TELEMETRY, POST_ATTRIBUTES)
%   2. Save Timeseries (write to ts_kv)
%   3. Route to Energy Billing Chain
%   4. Route to Alarm Processing Chain
% Figura: Screenshot Root Rule Chain

\subsubsection{Energy Billing Chain}
\label{sec:server-rule-chain-billing}

% TODO: NUEVO contenido (700 palabras)
% - Input: energy_kwh telemetry
% - Processing:
%   1. Query Asset attribute: tariff_usd_kwh
%   2. Script node (JS): cost_usd = energy_kwh × tariff
%   3. Save calculated attribute: monthly_cost_usd
%   4. Check threshold: > $100 → Notification
% - Output: Billing dashboard update
% Código: JavaScript node billing calculation

\subsubsection{Alarm Processing Chain}
\label{sec:server-rule-chain-alarms}

% TODO: NUEVO contenido (600 palabras)
% - Alarm rules:
%   * voltage_v < 207V OR > 253V → CRITICAL "Voltage Anomaly"
%   * current_a > 25A → MAJOR "Over-Current"
%   * temperature_c > 60°C → MINOR "High Temperature"
% - Actions:
%   * Create alarm entity
%   * Send notification (email, Slack webhook)
%   * Log to alarm_history table
% Tabla: Alarm types y severidades

\subsection{Gestión de Alarmas}
\label{sec:server-alarmas}

% TODO: MIGRAR y EXPANDIR (500 palabras)
% - Alarm lifecycle: NEW → ACTIVE → ACKNOWLEDGED → CLEARED
% - Alarm propagation: Device → Asset → Customer
% - Notification channels: Email, SMS, Webhook
% Figura: Flujo alarma end-to-end

\section{Diseño de Dashboards y Visualización}
\label{sec:server-dashboards}

% TODO: NUEVO contenido (1200 palabras)
% - Dashboards implementados:
%   1. Real-time Monitoring (telemetry widgets)
%   2. Historical Analytics (time-series charts)
%   3. Alarm Management (alarm table + map)
%   4. Energy Billing (bar charts, monthly aggregates)
% Longitud estimada: 4-5 páginas

\subsection{Dashboard de Monitoreo en Tiempo Real}
\label{sec:server-dashboard-realtime}

% TODO: NUEVO contenido (600 palabras)
% - Widgets:
%   * Gauge: voltage_v (current value)
%   * Line chart: power_w (last 1 hour)
%   * Map: Device locations with status color
%   * Alarm table: Active alarms (auto-refresh 5s)
% Figura: Screenshot dashboard monitoring

\subsection{Dashboard de Analítica Histórica}
\label{sec:server-dashboard-analytics}

% TODO: NUEVO contenido (600 palabras)
% - Widgets:
%   * Time-series chart: energy_kwh (last 30 days)
%   * Bar chart: Daily consumption per building
%   * Heatmap: Consumption by hour of day
%   * Statistics: Total kWh, Avg power factor
% Figura: Screenshot dashboard analytics

\section{APIs Programáticas y Webhooks}
\label{sec:server-apis}

% TODO: NUEVO contenido (800 palabras)
% - REST API ThingsBoard:
%   * Authentication: JWT tokens
%   * Device API: POST /api/device (provisioning)
%   * Telemetry API: GET /api/plugins/telemetry/DEVICE/{id}/values/timeseries
%   * Alarm API: GET /api/alarm/{entityType}/{entityId}
% - Webhooks configurados:
%   * Billing System: POST daily consumption → CRM
%   * Slack: POST critical alarms → #ami-alerts channel
% Código: cURL examples API calls

\section{Conclusiones del Capítulo}
\label{sec:server-conclusiones}

% TODO: Síntesis del capítulo (300-400 palabras)
% - Servidor ThingsBoard + Kafka implementado exitosamente
% - Kafka: 10K msgs/sec throughput validado
% - PostgreSQL: Query time < 100ms @ 1M registros
% - Rule Chains: Lógica de negocio compleja (billing, alarmas)
% - Siguiente capítulo: Validación end-to-end del sistema completo

